\begin{proof}[Proof of \cref{obj:lem:cosine-prior}]
\begin{enumerate}
\item \emph{Centering and square integrability.}
Fix the stated parameters and a coefficient-sign array \(\xi\). Write
\[
\begin{aligned}
\mathcal I
&=\{(j,\ell,a,b):1\leq j<\ell\leq d,\ 1\leq a,b\leq L\},\\
B&=\binom d2,\qquad M=BL^2,\qquad
\gamma_{\mathrm{prior}}=\frac{1}{C_0L^s\sqrt M},\\
V_{j\ell ab}(x)
&=2\cos(\pi a x_j)\cos(\pi b x_\ell).
\end{aligned}
\]
Choosing an index amounts to choosing a coordinate pair and two frequencies, so
\[
|\mathcal I|=\binom d2\,L^2=M.
\]
Since \(d\geq2\) and \(L\geq1\), we have \(B\geq1\), \(M>0\), and
\(\gamma_{\mathrm{prior}}>0\).

For every positive integer frequency \(a\),
\[
\int_0^1\cos(\pi a t)\,dt
=\frac{\sin(\pi a)}{\pi a}=0.
\]
Product integration therefore gives, for every \((j,\ell,a,b)\in\mathcal I\),
\[
\int_{[0,1]^d}V_{j\ell ab}(x)\,dx
=2\left(\int_0^1\cos(\pi a t)\,dt\right)
  \left(\int_0^1\cos(\pi b t)\,dt\right)
=0.
\]
The function \(m_\xi\) is a finite linear combination of continuous bounded
features, hence is Borel measurable and square integrable. Integrating its
finite sum yields
\[
\int_{[0,1]^d}m_\xi(x)\,dx
=\gamma_{\mathrm{prior}}
  \sum_{\alpha\in\mathcal I}\xi_\alpha
  \int_{[0,1]^d}V_\alpha(x)\,dx
=0.
\]
% lean: CausalSmith.Experimentation.BivariateSobolevDesignCapacity.mXi_mean_zero

\item \emph{Canonical components and reconstruction.}
Use the same finite cosine expression to evaluate \(m_\xi\) on
\(\mathbb R^d\). For \(\tau_j,\tau_\ell\in\mathbb R\), the canonical marginal
formulas are
\[
\begin{aligned}
g_j(m_\xi)(\tau_j)
&=\int_{[0,1]^d}
  \left.m_\xi(x)\right|_{x_j=\tau_j}\,dx,\\
g_{j\ell}(m_\xi)(\tau_j,\tau_\ell)
&=\int_{[0,1]^d}
  \left.m_\xi(x)\right|_{x_j=\tau_j,\ x_\ell=\tau_\ell}\,dx
  -g_j(m_\xi)(\tau_j)-g_\ell(m_\xi)(\tau_\ell),
  \qquad j<\ell.
\end{aligned}
\]
Here substitution fixes the indicated coordinates before integration; the
corresponding unused integration variables contribute factors of one.

For a feature indexed by
\(\alpha=(j_\alpha,\ell_\alpha,a_\alpha,b_\alpha)\), fixing coordinate \(j\)
leaves a centered cosine factor whenever \(j=j_\alpha\) or
\(j=\ell_\alpha\). If \(j\) equals neither coordinate, the integral is the
feature's mean. Thus, in all three cases,
\[
g_j(V_\alpha)(\tau_j)=0,
\qquad
g_j(m_\xi)(\tau_j)=0.
\]
After fixing \(j<\ell\), a feature survives integration precisely when its
coordinate pair is \((j,\ell)\). In every other configuration, at least one
of its two positive-frequency factors remains integrated and has mean
zero. The ordering of both pairs excludes a reversed match. Consequently,
\[
g_{j\ell}(V_\alpha)(\tau_j,\tau_\ell)
=
\begin{cases}
2\cos(\pi a_\alpha\tau_j)\cos(\pi b_\alpha\tau_\ell),
  &(j_\alpha,\ell_\alpha)=(j,\ell),\\
0,&(j_\alpha,\ell_\alpha)\ne(j,\ell),
\end{cases}
\]
and finite-sum linearity gives
\[
g_{j\ell}(m_\xi)(\tau_j,\tau_\ell)
=\gamma_{\mathrm{prior}}
  \sum_{a,b=1}^L\xi_{j\ell ab}
  \,2\cos(\pi a\tau_j)\cos(\pi b\tau_\ell).
\]
Summing these pair effects counts each feature exactly once:
\[
\sum_{j=1}^d g_j(m_\xi)(x_j)
+\sum_{j<\ell}g_{j\ell}(m_\xi)(x_j,x_\ell)
=\gamma_{\mathrm{prior}}
  \sum_{\alpha\in\mathcal I}\xi_\alpha V_\alpha(x)
=m_\xi(x).
\]
This proves the decomposition required in \cref{obj:ass:order-two},
pointwise and hence almost everywhere.
% lean: CausalSmith.Experimentation.BivariateSobolevDesignCapacity.mXi_orderTwo

\item \emph{A cutoff extension and its weak derivatives.}
To establish the component restriction bound, take an arbitrary real array
\[
c^{\mathrm{loc}}\in\mathbb R^{\{1,\ldots,L\}^2}
\]
and define, for
\(u_{\mathrm{loc}}=(u_1^{\mathrm{loc}},u_2^{\mathrm{loc}})\in\mathbb R^2\),
\[
\begin{aligned}
U_{\mathrm{loc}}(u_{\mathrm{loc}})
&=\sum_{a,b=1}^L c^{\mathrm{loc}}_{ab}
  \,2\cos(\pi a u_1^{\mathrm{loc}})
     \cos(\pi b u_2^{\mathrm{loc}}),\\
E_{\mathrm{coef}}
&=\sum_{a,b=1}^L(c^{\mathrm{loc}}_{ab})^2.
\end{aligned}
\]
Define the scalar cutoff and a representative of its almost-everywhere
derivative on the whole real line by
\[
q_{\mathrm{cut}}(t)=
\begin{cases}
0,&t<-1,\\
t+1,&-1\leq t<0,\\
1,&0\leq t\leq1,\\
2-t,&1<t\leq2,\\
0,&t>2,
\end{cases}
\qquad
q'_{\mathrm{cut}}(t)=
\begin{cases}
1,&-1<t<0,\\
-1,&1<t<2,\\
0,&\text{otherwise}.
\end{cases}
\]
In particular, the chosen derivative representative is zero at all four
joins. Put
\[
\begin{aligned}
\chi_{\mathrm{cut}}(u_{\mathrm{loc}})
&=q_{\mathrm{cut}}(u_1^{\mathrm{loc}})
  q_{\mathrm{cut}}(u_2^{\mathrm{loc}}),\\
\chi_{\mathrm{cut},1}(u_{\mathrm{loc}})
&=q'_{\mathrm{cut}}(u_1^{\mathrm{loc}})
  q_{\mathrm{cut}}(u_2^{\mathrm{loc}}),\\
\chi_{\mathrm{cut},2}(u_{\mathrm{loc}})
&=q_{\mathrm{cut}}(u_1^{\mathrm{loc}})
  q'_{\mathrm{cut}}(u_2^{\mathrm{loc}}),\\
G_{\mathrm{cut}}&=\chi_{\mathrm{cut}}U_{\mathrm{loc}},\\
D_\rho
&=\chi_{\mathrm{cut},\rho}U_{\mathrm{loc}}
  +\chi_{\mathrm{cut}}\partial_\rho U_{\mathrm{loc}},
  \qquad \rho\in\{1,2\}.
\end{aligned}
\]
These formulas give
\[
0\leq\chi_{\mathrm{cut}}\leq1,
\qquad
\chi_{\mathrm{cut},1}^2+\chi_{\mathrm{cut},2}^2\leq2.
\]
The extension \(G_{\mathrm{cut}}\) equals \(U_{\mathrm{loc}}\) on
\([0,1]^2\), is continuous, and is supported in \([-1,2]^2\).
Both \(G_{\mathrm{cut}}\) and \(D_\rho\) belong to
\(L^1(\mathbb R^2)\cap L^2(\mathbb R^2)\), since they are bounded and
compactly supported.

Fix the other coordinate and let \(\varphi\in C^1(\mathbb R;\mathbb R)\).
On each of the intervals \((-1,0)\), \((0,1)\), and \((1,2)\), ordinary
integration by parts applies to the corresponding slice of
\(G_{\mathrm{cut}}\), whose derivative there is the slice of \(D_\rho\).
Adding the three identities cancels the contributions at \(0\) and \(1\)
by continuity. The contributions at \(-1\) and \(2\) vanish because the
cutoff is zero. Extending the integrals to the real line therefore gives
\[
\int_{\mathbb R}
 \left.G_{\mathrm{cut}}(u_{\mathrm{loc}})
 \right|_{u_\rho^{\mathrm{loc}}=t}
 \varphi'(t)\,dt
=
-\int_{\mathbb R}
 \left.D_\rho(u_{\mathrm{loc}})
 \right|_{u_\rho^{\mathrm{loc}}=t}
 \varphi(t)\,dt.
\]
This holds for each \(\rho\in\{1,2\}\); hence \(D_\rho\) is the weak
coordinate derivative of \(G_{\mathrm{cut}}\).
% lean: CausalSmith.Experimentation.BivariateSobolevDesignCapacity.pairCosineExtensionPartial_weak_slice

\item \emph{Spatial energy on the nine support cells.}
For an integer \(\kappa\), integration of the cosine primitive gives
\[
\int_0^1\cos(\pi\kappa t)\,dt
=\mathbf1_{\{\kappa=0\}},
\qquad
\int_{-1}^2\cos(\pi\kappa t)\,dt
=3\mathbf1_{\{\kappa=0\}}.
\]
Indeed, for \(\kappa=0\) these are the interval lengths; for
\(\kappa\ne0\), all endpoint sine values vanish. For positive integer
frequencies \(a,b\), the identities
\[
\begin{aligned}
\cos(\pi a t)\cos(\pi b t)
&=\frac12\bigl(\cos(\pi(a-b)t)+\cos(\pi(a+b)t)\bigr),\\
\sin(\pi a t)\sin(\pi b t)
&=\frac12\bigl(\cos(\pi(a-b)t)-\cos(\pi(a+b)t)\bigr)
\end{aligned}
\]
therefore imply
\[
\begin{aligned}
\int_0^1\cos(\pi a t)\cos(\pi b t)\,dt
&=\int_0^1\sin(\pi a t)\sin(\pi b t)\,dt
 =\frac12\mathbf1_{\{a=b\}},\\
\int_{-1}^2\cos(\pi a t)\cos(\pi b t)\,dt
&=\int_{-1}^2\sin(\pi a t)\sin(\pi b t)\,dt
 =\frac32\mathbf1_{\{a=b\}}.
\end{aligned}
\]
Expanding the finite squares and using product integration eliminates
every off-diagonal frequency term. On the support square this yields
\[
\int_{[-1,2]^2}|U_{\mathrm{loc}}|^2
=4\left(\frac32\right)^2 E_{\mathrm{coef}}
=9E_{\mathrm{coef}}.
\]
Differentiating the polynomial gives
\[
\begin{aligned}
\partial_1U_{\mathrm{loc}}(u_{\mathrm{loc}})
&=-2\pi\sum_{a,b=1}^L
 a\,c^{\mathrm{loc}}_{ab}
 \sin(\pi a u_1^{\mathrm{loc}})
 \cos(\pi b u_2^{\mathrm{loc}}),\\
\partial_2U_{\mathrm{loc}}(u_{\mathrm{loc}})
&=-2\pi\sum_{a,b=1}^L
 b\,c^{\mathrm{loc}}_{ab}
 \cos(\pi a u_1^{\mathrm{loc}})
 \sin(\pi b u_2^{\mathrm{loc}}).
\end{aligned}
\]
The unit-square derivative energies satisfy
\[
\begin{aligned}
\int_{[0,1]^2}|\partial_1U_{\mathrm{loc}}|^2
&=\pi^2\sum_{a,b=1}^L a^2(c^{\mathrm{loc}}_{ab})^2
 \leq\pi^2L^2E_{\mathrm{coef}},\\
\int_{[0,1]^2}|\partial_2U_{\mathrm{loc}}|^2
&=\pi^2\sum_{a,b=1}^L b^2(c^{\mathrm{loc}}_{ab})^2
 \leq\pi^2L^2E_{\mathrm{coef}}.
\end{aligned}
\]
On \([-1,2]^2\), the corresponding energies are nine times these
frequency sums. In particular,
\[
\begin{aligned}
\int_{[-1,2]^2}|\partial_1U_{\mathrm{loc}}|^2
&=9\pi^2\sum_{a,b=1}^L a^2(c^{\mathrm{loc}}_{ab})^2
 \leq9\pi^2L^2E_{\mathrm{coef}},\\
\int_{[-1,2]^2}|\partial_2U_{\mathrm{loc}}|^2
&=9\pi^2\sum_{a,b=1}^L b^2(c^{\mathrm{loc}}_{ab})^2
 \leq9\pi^2L^2E_{\mathrm{coef}},\\
\int_{[-1,2]^2}\|\nabla U_{\mathrm{loc}}\|_2^2
&\leq18\pi^2L^2E_{\mathrm{coef}}.
\end{aligned}
\]
Since \(0\leq\chi_{\mathrm{cut}}\leq1\),
\[
\|G_{\mathrm{cut}}\|_{L^2(\mathbb R^2)}^2
\leq9E_{\mathrm{coef}}.
\]
For each coordinate, apply
\((r_1+r_2)^2\leq2r_1^2+2r_2^2\), valid for real \(r_1,r_2\), to the
two terms defining \(D_\rho\). Summing the two coordinate inequalities
and using the cutoff derivative bound gives
\[
\frac12(D_1^2+D_2^2)
\leq
\chi_{\mathrm{cut}}^2
 \bigl(|\partial_1U_{\mathrm{loc}}|^2+
       |\partial_2U_{\mathrm{loc}}|^2\bigr)
+2|U_{\mathrm{loc}}|^2.
\]
Thus
\[
D_1^2+D_2^2
\leq2\|\nabla U_{\mathrm{loc}}\|_2^2+4|U_{\mathrm{loc}}|^2
\]
on the support square. Integration proves
\[
\|D_1\|_{L^2(\mathbb R^2)}^2+
\|D_2\|_{L^2(\mathbb R^2)}^2
\leq(36+36\pi^2L^2)E_{\mathrm{coef}}.
\]
% lean: CausalSmith.Experimentation.BivariateSobolevDesignCapacity.pairCosineExtensionPartial_spatial_energy_le

\item \emph{Fourier endpoint bounds.}
For an integrable function
\(\widetilde G:\mathbb R^2\to\mathbb C\), use the unitary angular-frequency
Fourier transform
\[
(\mathcal F_2\widetilde G)(\omega_{\mathrm{loc}})
=\frac{1}{2\pi}\int_{\mathbb R^2}
 e^{-\mathrm i\langle\omega_{\mathrm{loc}},u_{\mathrm{loc}}\rangle}
 \widetilde G(u_{\mathrm{loc}})\,du_{\mathrm{loc}},
\qquad \omega_{\mathrm{loc}}\in\mathbb R^2.
\]
For \(\lambda\in[0,1]\) and
\(\widetilde G\in L^1(\mathbb R^2)\cap L^2(\mathbb R^2)\), define
\[
\mathcal E_\lambda(\widetilde G)
=\int_{\mathbb R^2}
 |\mathcal F_2\widetilde G(\omega_{\mathrm{loc}})|^2
 \left(1+\frac{\|\omega_{\mathrm{loc}}\|_2^2}{2}\right)^\lambda
 \,d\omega_{\mathrm{loc}}.
\]
The slice integration-by-parts identities, applied to the real and
imaginary parts of the Fourier exponential, give
\[
\mathcal F_2D_\rho(\omega_{\mathrm{loc}})
=\mathrm i\,\omega_\rho^{\mathrm{loc}}\,
 \mathcal F_2G_{\mathrm{cut}}(\omega_{\mathrm{loc}}),
\qquad \rho\in\{1,2\}.
\]
Plancherel's identity consequently yields
\[
\begin{aligned}
\mathcal E_0(G_{\mathrm{cut}})
&=\|G_{\mathrm{cut}}\|_2^2
 \leq9E_{\mathrm{coef}}
 \leq16E_{\mathrm{coef}},\\
\mathcal E_1(G_{\mathrm{cut}})
&=\|G_{\mathrm{cut}}\|_2^2
 +\frac12\bigl(\|D_1\|_2^2+\|D_2\|_2^2\bigr)\\
&\leq(27+18\pi^2L^2)E_{\mathrm{coef}}.
\end{aligned}
\]
The calibration uses \(L^2\geq1\):
\[
27+18\pi^2L^2
\leq(48+32\pi^2)L^2=C_0^2L^2.
\]
Hence the required endpoint estimates are
\[
\mathcal E_0(G_{\mathrm{cut}})\leq16E_{\mathrm{coef}},
\qquad
\mathcal E_1(G_{\mathrm{cut}})\leq C_0^2L^2E_{\mathrm{coef}}.
\]
% lean: CausalSmith.Experimentation.BivariateSobolevDesignCapacity.pairCosineExtension_fourier_endpoints

\item \emph{Fractional energy and the restriction norm.}
Define the nonnegative energy densities on \(\mathbb R^2\) by
\[
\begin{aligned}
\eta_0(\omega_{\mathrm{loc}})
&=|\mathcal F_2G_{\mathrm{cut}}(\omega_{\mathrm{loc}})|^2,\\
\eta_1(\omega_{\mathrm{loc}})
&=\eta_0(\omega_{\mathrm{loc}})
 \left(1+\frac{\|\omega_{\mathrm{loc}}\|_2^2}{2}\right).
\end{aligned}
\]
With the convention \(r^0=1\) for \(r\geq0\), their pointwise factorization is
\[
\eta_0^{\,1-s}\eta_1^{\,s}
=\eta_0
 \left(1+\frac{\|\omega_{\mathrm{loc}}\|_2^2}{2}\right)^s.
\]
The integral form of Hölder's inequality with nonnegative weights
\(1-s\) and \(s\), whose sum is one, therefore gives
\[
\mathcal E_s(G_{\mathrm{cut}})
\leq
\mathcal E_0(G_{\mathrm{cut}})^{1-s}
\mathcal E_1(G_{\mathrm{cut}})^s.
\]
This form includes endpoint weights and zero energy. Since
\(16\leq C_0^2\), the endpoint bounds imply
\[
\begin{aligned}
\mathcal E_s(G_{\mathrm{cut}})
&\leq
 (C_0^2E_{\mathrm{coef}})^{1-s}
 (C_0^2E_{\mathrm{coef}}L^2)^s\\
&=C_0^2L^{2s}E_{\mathrm{coef}}.
\end{aligned}
\]
In particular, when \(E_{\mathrm{coef}}=0\), every coefficient and the
extension vanish, and this bound is zero.

The squared restriction norm is the extension infimum
\[
N_{2,s}(U_{\mathrm{loc}}|_{[0,1]^2})^2
=
\inf_{\substack{
 \widetilde G\in L^1(\mathbb R^2;\mathbb C)\cap L^2(\mathbb R^2;\mathbb C)\\
 \widetilde G=U_{\mathrm{loc}}\ \text{a.e.\ on }[0,1]^2}}
 \mathcal E_s(\widetilde G).
\]
The continuous compactly supported extension \(G_{\mathrm{cut}}\) is
admissible in this infimum. We obtain
\[
N_{2,s}(U_{\mathrm{loc}}|_{[0,1]^2})^2
\leq
\mathcal E_s(G_{\mathrm{cut}})
\leq C_0^2L^{2s}
 \sum_{a,b=1}^L(c^{\mathrm{loc}}_{ab})^2.
\]
% lean: CausalSmith.Experimentation.BivariateSobolevDesignCapacity.pairCosinePolynomial_restriction_bound

\item \emph{Component budget estimates.}
For each \(j<\ell\), define the normalized pair coefficients and their
energy by
\[
c^{\mathrm{prior}}_{j\ell ab}
=\gamma_{\mathrm{prior}}\xi_{j\ell ab},
\qquad
E^{\mathrm{prior}}_{j\ell}
=\sum_{a,b=1}^L(c^{\mathrm{prior}}_{j\ell ab})^2,
\qquad 1\leq a,b\leq L.
\]
Every sign has square one, so
\[
E^{\mathrm{prior}}_{j\ell}
=\sum_{a,b=1}^L\gamma_{\mathrm{prior}}^2
=L^2\gamma_{\mathrm{prior}}^2.
\]
The canonical pair formula and the restriction bound now give
\[
N_{2,s}(g_{j\ell}(m_\xi))^2
\leq C_0^2L^{2s}E^{\mathrm{prior}}_{j\ell}
=C_0^2L^{2s}L^2\gamma_{\mathrm{prior}}^2.
\]
The main effects are zero; the zero extension certifies their zero
restriction norms. Hence
\[
\sum_{j=1}^dN_{1,s}(g_j(m_\xi))^2=0.
\]
Summing the common pair bound over the \(B\) coordinate pairs yields
\[
\sum_{j=1}^dN_{1,s}(g_j(m_\xi))^2
+\sum_{j<\ell}N_{2,s}(g_{j\ell}(m_\xi))^2
\leq B\,C_0^2L^{2s}L^2\gamma_{\mathrm{prior}}^2.
\]
% lean: CausalSmith.Experimentation.BivariateSobolevDesignCapacity.mXi_pooledBudget

\item \emph{Cube inner products.}
Take
\[
\alpha=(j,\ell,a,b)\in\mathcal I,
\qquad
\beta=(j',\ell',a',b')\in\mathcal I.
\]
Product integration factors the feature inner product as
\[
\begin{aligned}
\int_{[0,1]^d}V_\alpha(x)V_\beta(x)\,dx
=4\prod_{\iota=1}^d\int_0^1
&\cos(\pi a t)^{\mathbf1_{\{\iota=j\}}}
 \cos(\pi b t)^{\mathbf1_{\{\iota=\ell\}}}\\
&{}\cdot
 \cos(\pi a't)^{\mathbf1_{\{\iota=j'\}}}
 \cos(\pi b't)^{\mathbf1_{\{\iota=\ell'\}}}\,dt .
\end{aligned}
\]
First suppose \(j=j'\). If also \(\ell=\ell'\), the unit-interval cosine
integrals give
\[
\int_{[0,1]^d}V_\alpha V_\beta
=4\left(\frac12\mathbf1_{\{a=a'\}}\right)
  \left(\frac12\mathbf1_{\{b=b'\}}\right)
=\mathbf1_{\{\alpha=\beta\}}.
\]
If instead \(\ell\ne\ell'\), then \(\ell\) is distinct from both
\(j'=j\) and \(\ell'\). Its factor occurs exactly once and has integral
zero.

Now suppose \(j\ne j'\). If \(j=\ell'\), the ordering gives
\[
j'<\ell'=j<\ell,
\]
so \(\ell\) is distinct from both coordinates of \(\beta\) and occurs
exactly once. If \(j\ne\ell'\), coordinate \(j\) is distinct from both
coordinates of \(\beta\) and occurs exactly once. In each of these
unmatched-coordinate cases, the product factorization contains a
positive-frequency cosine mean and consequently gives
\[
\int_{[0,1]^d}V_\alpha(x)V_\beta(x)\,dx=0.
\]
Combining the cases proves
\[
\int_{[0,1]^d}V_\alpha(x)V_\beta(x)\,dx
=\mathbf1_{\{\alpha=\beta\}},
\qquad \alpha,\beta\in\mathcal I.
\]
% lean: CausalSmith.Experimentation.BivariateSobolevDesignCapacity.pairFeature_cube_inner_product

\item \emph{Normalization and class membership.}
Expanding the finite square and applying the cube inner-product identity
gives
\[
\begin{aligned}
\int_{[0,1]^d}m_\xi(x)^2\,dx
&=\gamma_{\mathrm{prior}}^2
 \sum_{\alpha,\beta\in\mathcal I}
 \xi_\alpha\xi_\beta\mathbf1_{\{\alpha=\beta\}}\\
&=\gamma_{\mathrm{prior}}^2
 \sum_{\alpha\in\mathcal I}\xi_\alpha^2
=\gamma_{\mathrm{prior}}^2M
=\frac{1}{C_0^2L^{2s}}.
\end{aligned}
\]
Thus the normalized coefficient energy satisfies
\[
C_0^2L^{2s}
 \sum_{\alpha\in\mathcal I}
  (\gamma_{\mathrm{prior}}\xi_\alpha)^2
=C_0^2L^{2s}\gamma_{\mathrm{prior}}^2M
=1.
\]
Since \(M=BL^2\), this also calibrates the common pair budget:
\[
B\,C_0^2L^{2s}L^2\gamma_{\mathrm{prior}}^2
=C_0^2L^{2s}\gamma_{\mathrm{prior}}^2M
=1.
\]
The component estimates therefore establish
\[
\sum_{j=1}^dN_{1,s}(g_j(m_\xi))^2
+\sum_{j<\ell}N_{2,s}(g_{j\ell}(m_\xi))^2
\leq1,
\]
as required in \cref{obj:ass:pooled-budget}. Together with measurability,
square integrability, centering, and the canonical order-two
reconstruction, this proves
\[
m_\xi\in\mathcal H_{d,s}
\]
by \cref{obj:def:sobolev-class}.
% lean: CausalSmith.Experimentation.BivariateSobolevDesignCapacity.mXi_pooledBudget

\item \emph{Transfer to the first sampled covariate.}
Let \(n\geq1\) and let \(X\) satisfy \cref{obj:ass:uniform-draw}.
Taking the first coordinate marginal of its product sampling law gives
\[
\operatorname{Law}(X_1)=\operatorname{Unif}([0,1]^d).
\]
The feature products are measurable and bounded, so integration against
this marginal yields
\[
\mathbb E[V_\alpha(X_1)V_\beta(X_1)]
=\int_{[0,1]^d}V_\alpha(x)V_\beta(x)\,dx
=\mathbf1_{\{\alpha=\beta\}},
\qquad \alpha,\beta\in\mathcal I.
\]
There are \(M\) feature coordinates. Taking these expectations entrywise
therefore proves
\[
\mathbb E[V(X_1)V(X_1)^{\mathsf T}]=I_M.
\]
The unmatched-coordinate cases include distinct pairs with one shared
original coordinate, so the identity covers those pairs as well.
% lean: CausalSmith.Experimentation.BivariateSobolevDesignCapacity.uniformDraw_pairFeature_inner_product
\end{enumerate}
\end{proof}
